\documentclass[main.tex]{subfiles}


\begin{document}

%from pagenumber to up
% \phantomsection 
% \hypertarget{toc}{}
 

 

% номер секции вручную
\setcounter{section}{31
}
\addtocounter{section}{-1} 

\section{Опыт Эрстеда}

% Постоянный магнит~--- не единственный источник магнитного поля. 
Магнитными свойствами обладает также электрический ток. Обнаружить на опыте эти свойства тока и изучить их удалось в опыте Эрстеда.

Схема опыта Эрстеда показана на рис.\,\ref{fig:p145}.

    \begin{figure}[H]
    \centering
\begin{circuitikz}[circuit ee IEC,font=\small,thin,
    line cap=round,line join=round,
>={Stealth[inset=0pt,round,length=3mm, angle'=20]},
vec/.style={>={Stealth[inset=0pt,round,length=3.5mm, angle'=20]}}]

    \ctikzset{bipoles/americaninductor/coils=3}


\path[gray,rounded corners=1pt] (.75,1.9) --++ (0,.2) --++ (-.25,0) coordinate (a) --++ (-.25,0) --++ (0,-.2) -- cycle; 

\path[gray,rounded corners=1pt] (.85,.1) --++ (0,-.2) --++ (-.35,0) coordinate (b) --++ (-.35,0) --++ (0,.2) -- cycle; 


%%%%%%%%%%%%%%%%%%%%%%%%%%%%%%%%
%%%%%%%%%%%%%%%%%%%%%%%%%%%%%%%%
% arrow
\def\dc{.2}


\path[rotate around={40:(5.5,1)}] (5.5,1) ++ (0,\dc) coordinate (ac) --++ ({3*\dc},{-\dc}) coordinate (rc) --++ ({-3*\dc},{-\dc}) coordinate (bc) --++ ({-3*\dc},{\dc}) coordinate (lc) --++ ({3*\dc},{\dc});


\begin{scope}[even odd rule]
\clip [] (lc) -- (ac) -- (rc) -- (bc) -- cycle (5.5,1) circle (1);     
\fill[lightgray,semitransparent] (5.5,1) circle (1);
\end{scope}
\draw (5.5,1) circle (1);
\draw[shift={(5.5,1)}] (0,1) --++ (0,-.2) (1,0) node[right] {М} --++ (-.2,0) (0,-1) --++ (0,.2) (-1,0) --++ (.2,0);


%nee
\fill[red,semitransparent] (lc) -- (ac) -- (bc) -- cycle;
\fill[NavyBlue,semitransparent] (rc) -- (ac) -- (bc) -- cycle;
\draw (ac) -- (rc) -- (bc) -- (lc) -- cycle;

\node at (5.5,1) [circle,fill=,inner sep=0pt,minimum size=1mm,label=below:] {};


%cond
\draw[] (a) --++ (0,1) 
to[make contact={info={[black]$K$}}]++ 
(5,0) node[above,black] {$O$} coordinate (ar);
\draw[] (b) --++ (0,-1) --++ (5,0) node[below,black] {$O'$} coordinate (br);
\draw[thick,brown] (br) -- (ar);
% to[normal open switch,l=$K$,black]++ 


%terminals
\begin{scope}
    \clip[] (0,2) rectangle++ (1,.2);
    \filldraw[gray,rounded corners=1pt] (.75,1.9) --++ (0,.2) --++ (-.25,0) coordinate (a) --++ (-.25,0) --++ (0,-.2) -- cycle; 
\end{scope}

\begin{scope}
    \clip[] (0,0) rectangle++ (1,-.2);
    \filldraw[gray,rounded corners=1pt] (.85,.1) --++ (0,-.2) --++ (-.35,0) coordinate (b) --++ (-.35,0) --++ (0,.2) -- cycle; 
\end{scope}


%batt
\begin{scope}
    \clip[rounded corners=1mm] (0,0) rectangle++ (1,2);
    \fill[gray,nearly transparent] (0,0) rectangle++ (1,1);
    \fill[yellow,nearly transparent] (0,1) rectangle++ (1,1);
\end{scope}
\draw[rounded corners=1mm] (0,0) rectangle++ (1,2);





%labels
\node[above] at (.5,0) {$-$};
\node[below] at (.5,2) {$+$};
\node[left] at (0,1) {Б}; 





\end{circuitikz}
\caption{Опыт Эрстеда: тока в цепи нет}
\label{fig:p145}
\end{figure}


Над магнитной стрелкой~М (стрелкой компаса) расположен провод~$OO'$, подключенный через ключ~$K$ к батарейке~Б. Если ключ разомкнут (в цепи из батарейки и проводников нет тока), то магнитная стрелка указывает на (географический) север Земли.

Пусть теперь ключ в рассмотренной схеме замыкают (рис.\,\ref{fig:p146}).


    \begin{figure}[H]
    \centering
\begin{circuitikz}[circuit ee IEC,font=\small,thin,
    line cap=round,line join=round,
>={Stealth[inset=0pt,round,length=3mm, angle'=20]},
vec/.style={>={Stealth[inset=0pt,round,length=3.5mm, angle'=20]}}]

    \ctikzset{bipoles/americaninductor/coils=3}


\path[gray,rounded corners=1pt] (.75,1.9) --++ (0,.2) --++ (-.25,0) coordinate (a) --++ (-.25,0) --++ (0,-.2) -- cycle; 

\path[gray,rounded corners=1pt] (.85,.1) --++ (0,-.2) --++ (-.35,0) coordinate (b) --++ (-.35,0) --++ (0,.2) -- cycle; 


%%%%%%%%%%%%%%%%%%%%%%%%%%%%%%%%
%%%%%%%%%%%%%%%%%%%%%%%%%%%%%%%%
% arrow
\def\dc{.2}


\path[rotate around={0:(5.5,1)}] (5.5,1) ++ (0,\dc) coordinate (ac) --++ ({3*\dc},{-\dc}) coordinate (rc) --++ ({-3*\dc},{-\dc}) coordinate (bc) --++ ({-3*\dc},{\dc}) coordinate (lc) --++ ({3*\dc},{\dc});


\begin{scope}[even odd rule]
\clip [] (lc) -- (ac) -- (rc) -- (bc) -- cycle (5.5,1) circle (1);     
\fill[lightgray,semitransparent] (5.5,1) circle (1);
\end{scope}
\draw (5.5,1) circle (1);
\draw[shift={(5.5,1)}] (0,1) --++ (0,-.2) (1,0) node[right] {М} --++ (-.2,0) (0,-1) --++ (0,.2) (-1,0) --++ (.2,0);


%nee
\fill[red,semitransparent] (lc) -- (ac) -- (bc) -- cycle;
\fill[NavyBlue,semitransparent] (rc) -- (ac) -- (bc) -- cycle;
\draw (ac) -- (rc) -- (bc) -- (lc) -- cycle;

\node at (5.5,1) [circle,fill=,inner sep=0pt,minimum size=1mm,label=below:] {};


%cond
\draw[] (a) --++ (0,1) 
% to[make contact={info={[black]$K$}}]++ 
--++ (2.5,0) node[black,above] {$K$} --++ (2.5,0) node[above,black] {$O$} coordinate (ar);
\draw[] (b) --++ (0,-1) --++ (5,0) node[below,black] {$O'$} coordinate (br);
\draw[brown,thick] (br) -- (ar);
% to[normal open switch,l=$K$,black]++ 


%terminals
\begin{scope}
    \clip[] (0,2) rectangle++ (1,.2);
    \filldraw[gray,rounded corners=1pt] (.75,1.9) --++ (0,.2) --++ (-.25,0) coordinate (a) --++ (-.25,0) --++ (0,-.2) -- cycle; 
\end{scope}

\begin{scope}
    \clip[] (0,0) rectangle++ (1,-.2);
    \filldraw[gray,rounded corners=1pt] (.85,.1) --++ (0,-.2) --++ (-.35,0) coordinate (b) --++ (-.35,0) --++ (0,.2) -- cycle; 
\end{scope}


%batt
\begin{scope}
    \clip[rounded corners=1mm] (0,0) rectangle++ (1,2);
    \fill[gray,nearly transparent] (0,0) rectangle++ (1,1);
    \fill[yellow,nearly transparent] (0,1) rectangle++ (1,1);
\end{scope}
\draw[rounded corners=1mm] (0,0) rectangle++ (1,2);





%labels
\node[above] at (.5,0) {$-$};
\node[below] at (.5,2) {$+$};
\node[left] at (0,1) {Б}; 





\end{circuitikz}
\caption{Опыт Эрстеда: ток в цепи есть}
\label{fig:p146}
\end{figure}


Замыкание ключа~$K$ приводит к протеканию тока в цепи из батарйки~Б и проводников. При достаточно большом токе в цепи магнитная стрелка~М устанавливается \emph{перпендикулярно проводу}~$OO'$. Оказывается, отклонение магнитной стрелки вызвано действием на нее со стороны провода~$OO'$ с током (стрелка и провод с током взаимодействуют друг с другом подобно тому, как магнитная стрелка взаимодействует с постоянным магнитом). При размыкании ключа магнитная стрелка возвращается в свое начальное положение.

Опыт Эрстеда показывает, что \emph{электрический ток создает в окружающем пространстве магнитное
поле.}



























 
\end{document}